\section{The remaining small-$r$ local chart of weak-C2 R8}
\label{sec:r17-r8-small-r}

\begin{theorem}
Retain the original physical stationary weak-C2 source, all strict gaps,
nonnegative event coefficients, the original positive-infinitesimal
total-scale intersection, and every inherited residual exclusion.
Suppose the source is in R8 and $r$ is infinitesimal. Thus $t=y/B$
and $g=G/B$ are positive bounded units, while $B,h=H/B,j,r$ are
positive infinitesimals. Then the entire original four-state parent is
a finite tensor SOS for every independent pair of infinitesimal or zero
testing half-angles $u,v$, with either sign and all real or complex states.
No source opening is compared with another opening or with a testing rate.
\end{theorem}

This is the missing local component of the R8 join. Its remote tests
and its $r$-unit local chart are supplied separately by the R17
joint-$D$ argument. The proof below does not rely on an unproved
extension of the RC normalization across $r=0$.

\paragraph{Source entry and the three independent hard scales.}
Write $a=z/B$, $l=1-h$, $x=aB$, and
\[
 j=1+xr-l/t,\qquad g=1+a(hB+r-1).
\]
The R8 and small-$r$ hypotheses imply that $a$ is bounded with residue
$\rho\in[0,1)$: $g_0=1-\rho>0$. Also $t_0=1$.
The unchanged exact source formulas WE1, jointly at $B=h=r=0,t=1$,
give
\[
 (nr)_0=(Ah)_0=(Dj)_0=\frac1{2(1+\rho)},\qquad
 \left(\frac{m}{a^2}\right)_0=\frac1{2(1-\rho^2)},
 \quad C_0=\left(\frac\alpha{B^2}\right)_0
       =\frac{1-\rho}{2(1+\rho)}>0.                 \tag{RD1}
\]
These are joint regular reductions. After replacing $n,A,D$ by
$nr,Ah,Dj$, all source denominators are products of units
$t,g,1-h,1-hB,\Delta$ with $\Delta_0=1+\rho$.
Thus $n,A,D$ are all positive unbounded, with completely unrestricted
relative rates, while $nr,Ah,Dj,\mu=m/a^2,C,\alpha_*=\alpha/B^2$
are positive bounded units. In particular $n=0$ is impossible on
this chart; no inverse of a vanishing event coefficient is being used.

\paragraph{A general exact centered-$D$ gauge of the complete WA matrix.}
Use the unnormalized complete WA3--WA8 identity, with the WE4 source
rescaling of the second scalar root:
\[
 \mathcal R=(\xi,\eta,R_1,R_2',T)^T,
 \quad \xi=X+ruZ,\quad\eta=Y-huW,
\]
\[
 R_1=uW-t^{-1}vZ,\quad R_2'=avW-uZ,\quad T=uvW-BZ,
 \qquad F/4=\mathcal R^TG\mathcal R.                \tag{RD2}
\]
Here $G$ contains the complete hard block, both complete first-normal
rows, and the complete scalar block. No packet or normal is omitted.
All its entries are affine in the formal event coefficient $D$.

For $U=u^2$, $V=v^2$, define three five-component columns
\begin{align*}
 d&=(uv,-1,l,ruv,-arv)^T,\\
 s&=(0,0,(x-U)v,-(B-V/t)u,U-aV/t)^T,\\
 k&=(-2rtu,0,0,-r^2tu,ar^2t)^T.
\end{align*}
The physical root relation and the centered $D$ event are exactly
\[
 s^T\mathcal R=0,\qquad
 L_D:=d^T\mathcal R
   =uW+uvX-Y-vZ+jvZ.                               \tag{RD3}
\]
The following is a matrix identity, valid before any residue operation:
\[
 \partial_DG-dd^T-\frac{sk^T+ks^T}{2}=jH_D,        \tag{RD4}
\]
where $H_D$ is symmetric and its only nonzero entries are
\[
 (H_D)_{13}=tuv,\quad (H_D)_{15}=-tu,\quad
 (H_D)_{34}=rtuv,\quad (H_D)_{35}=-artv.           \tag{RD5}
\]
Indices in RD5 refer to the displayed order in RD2 and start at one.
This explicitly resolves the failure of the ungauged matrix
$\partial_DG-dd^T$ to be divisible by $j$.

Let $G^{(0)}=G|_{D=0}$ denote a formal matrix substitution, with every
other source coefficient kept at its actual value. This is not a
physical modified source and no physicality theorem is applied to it.
Because the term involving $s$ vanishes exactly on the physical roots,
RD4 gives the entire-parent identity
\[
 F/4=D L_D^2+\mathcal R^T\{G^{(0)}+DjH_D\}\mathcal R.
                                                               \tag{RD6}
\]
The first-normal and hard-block entries in RD4 are part of this identity,
not merely its scalar restriction. No inverse of $u$, $v$, $j$, or a
state coordinate was used in deriving RD3--RD6.

\paragraph{The bounded $D$-free matrix through the $a$ opening.}
Normalize only the first two roots in $G^{(0)}$ initially, by
$\widehat\xi=\sqrt n\,\xi$ and $\widehat\eta=\sqrt A\,\eta$.
The resulting matrix, in
$(\widehat\xi,\widehat\eta,R_1,R_2',T)$, is a finite polynomial in
$u,v$ with bounded pure-source coefficients. Here is the complete
coefficient check needed for this assertion.

The hard entries after $D=0$ are
\begin{align*}
 h_{xx}/n&=(1+U)(1+V)+[m+UV/2+(1+V)/2]/n,\\
 h_{xy}/\sqrt{nA}&=-uv(m+1/2)/\sqrt{nA},\\
 h_{yy}/A&=(1+U)(1+V)
       +[C(1+V)+1/2+mUV+(\alpha+1/2)(1+U)V]/A.
\end{align*}
Every added coefficient is bounded and has zero residue.
In the first normal, the complete WA6 parameters now have the forms
\[
 p_0=t(nr+r/2),\quad p_1=t(nr+r/2),\quad
 q_1/a=nr+Bp_1,
\]
\[
 f_x=-mh+l/2-p_0-q_1,
 \quad (-m+q_1U)/a=-a\mu+(nr+Bp_1)U.
\]
They are bounded before multiplication by $1/\sqrt n$.
In the complete second normal, all WA7 coefficients involve only
the bounded products $Ah,m,C,\alpha,r,h$, and
\[
 q_y/a=-k_z-(Ah)/t-Bk_{uv}
\]
is bounded. Multiplication by $1/\sqrt A$ therefore bounds that
entire normal row as well. No ratio $n/A$, $A/n$, $h/r$, or $r/h$
appears.

For the scalar block use the exact WE5 expressions with $D=0$.
All occurrences of $n$ are through $nr^2=(nr)r$, and all occurrences
of $A$ are through $Ah^2=(Ah)h$. The only possible new $a$ denominator
is
\[
 Q_*=(b_2-\alpha_*)/a,\quad
 b_2=2a^2\mu h+C(1-h)^2+(Ah)h+(B^2\alpha_*+1/2)h^2. \tag{RD7}
\]
The exact WE1 expressions for $b_2$ and $\alpha_*$ involve no $D$
and no $j$ denominator. At $a=0$ their difference vanishes identically
for arbitrary $B,h,r,t$. Its cleared numerator is consequently
divisible by $a$; the canceled denominator at $B=h=r=0,t=1$ is
$2(a-1)(a+1)$, a unit for $0\le\rho<1$. The divided numerator
vanishes at that joint corner for every $a$. Hence $Q_*$ is bounded
(in fact infinitesimal), independently of $a,h,r,j$ rates.
This proves the bound on
\[
 f_s/a=\tfrac12[Q_*-a\,nr^2-aBt(n+1/2)r^2]
\]
and on the whole remaining scalar block. In particular the old
assumption that $j$ was a unit is unnecessary after the exact $D$
channel has been removed by RD6.

\paragraph{Simultaneous normalization of all three hard scales.}
Use the final five roots
\[
 \mathcal N=(\widehat\xi,\widehat\eta,\widehat D,R_2',T)^T,
 \qquad\widehat D=\sqrt D\,L_D.                    \tag{RD8}
\]
Solving only the source-unit coefficient $l=1-h$ in RD3 gives
\[
 R_1=\frac1l\left\{
 \frac{\widehat D}{\sqrt D}
 +\frac{\widehat\eta}{\sqrt A}
 -\frac{uv\widehat\xi}{\sqrt n}
 -ruvR_2'+arvT\right\}.                            \tag{RD9}
\]
Every coefficient in RD9 is bounded and separated. In particular,
independent small $r,h,j$ cause no ratio of their square roots.
Let $P(u,v)$ be the explicit matrix mapping $\mathcal N$ to
$\mathcal R$ according to RD9 and
$\xi=\widehat\xi/\sqrt n$, $\eta=\widehat\eta/\sqrt A$.
Then exactly
\[
 F/4=\mathcal N^T\widehat G\mathcal N,
 \qquad
 \widehat G=E_{33}+P^T\{G^{(0)}+DjH_D\}P.          \tag{RD10}
\]
The preceding bounds and RD9 prove that every coefficient of this
finite testing-polynomial matrix is bounded pure source. The factor
$Dj$ is a unit, while every entry of $H_D$ has positive testing degree.
Thus its contribution to the residue at $u=v=0$ is zero.

The normalized $n$ and $A$ hard diagonals have residue one. Since
$1/\sqrt n$, $1/\sqrt A$, $1/\sqrt D$, $r,h$ are infinitesimal,
the entire $R_1$ row of $P(0,0)$ has zero residue. The remaining
source scalar diagonals have residues $\mu_0$ and $\alpha_{*,0}=C_0$.
All cross entries have zero residue. Consequently
\[
 \overline{\widehat G}
   =\operatorname{diag}\left(1,1,1,
              \frac1{2(1-\rho^2)},
              \frac{1-\rho}{2(1+\rho)}\right)\succ0. \tag{RD11}
\]
This is the full five-root Gram residue, not a Hessian or state slice.

The positive pure-source matrix $\widehat G(0,0)$ has bounded inverse.
Whiten it purely in the source field and apply the finite separated
coefficient construction WA10 to the finitely many nonconstant
testing monomials. Their values are infinitesimal for arbitrary
independent infinitesimal $u,v$. This produces a finite SOS in
$K_c\otimes_{\mathbb R}K_t$ without any mixed inverse or mixed radical.
The argument includes either zero test and both signs because all
identities are polynomial in $u,v$. Applying it to real and imaginary
state parts proves the complex-state statement.

\paragraph{Verification boundary.}
The focused checker imports only the established complete six-contact
WA identity and WA10. It verifies the new entire $D$-matrix gauge,
the physical syzygy, the exact original-event centering, the entire
parent equality RD6, all fifteen normalized coefficient denominators
at the joint opening, the full residue RD11, and the exact source
quotient through $a=0$. Eighteen new exact checks pass. No numerical
positivity test or unproved scalar-to-parent lift is used.
